% Build:  tectonic paper/rusty-numbers-note.tex
\documentclass[10pt,reqno]{amsart}

\usepackage[T1]{fontenc}
\usepackage[utf8]{inputenc}
\usepackage{amsmath,amssymb,amsthm}
\usepackage{booktabs}
\usepackage{array}
\usepackage{microtype}
\usepackage[margin=0.8in]{geometry}
\usepackage[colorlinks=true,linkcolor=black,citecolor=black,urlcolor=blue]{hyperref}

\theoremstyle{plain}
\newtheorem{theorem}{Theorem}
\newtheorem{lemma}[theorem]{Lemma}
\newtheorem{corollary}[theorem]{Corollary}
\newtheorem{proposition}[theorem]{Proposition}
\newtheorem{conjecture}[theorem]{Conjecture}

\theoremstyle{definition}
\newtheorem{definition}[theorem]{Definition}
\newtheorem{observation}[theorem]{Observation}

\theoremstyle{remark}
\newtheorem{remark}[theorem]{Remark}

\numberwithin{theorem}{section}

% compact vertical spacing
\setlength{\abovedisplayskip}{5pt}
\setlength{\belowdisplayskip}{5pt}
\setlength{\abovedisplayshortskip}{3pt}
\setlength{\belowdisplayshortskip}{3pt}
\makeatletter
\def\thm@space@setup{\thm@preskip=5pt \thm@postskip=5pt}
\makeatother

\newcommand{\N}{\mathbb{N}}
\DeclareMathOperator{\mex}{mex}

\begin{document}

\title[Machine verification for the rusty numbers]{A machine-verified family of anti-recurrence
sequences: the rusty numbers and 248 further linear forms}

\author{Andrew Hingston}
\email{\texttt{hingstonandrew@icloud.com}}

\date{18 August 2026}

\subjclass[2020]{68Q45, 68R15, 11B85}
\keywords{anti-recurrence sequence, automatic sequence, DFAO, decision procedure, mex}

\begin{abstract}
Fokkink and Joshi's anti-recurrence conjecture (arXiv:2506.13337; published as
\emph{Anti-recurrence sequences}, Integers \textbf{26} (2026), \#A24, Conjecture~2) predicts
that for every positive integral linear form $a=(a_1,\dots,a_k)$ of dimension $k>1$, the
anti-recurrence sequence $A_n$ satisfies $A_n-\kappa n$ $\tau$-automatic, $\kappa=k\tau+1$,
$\tau$ the trace. The authors prove this, their Theorem~4, only under an
``$A_1$-boundedness'' restriction, and single out the \emph{rusty numbers} $a=(1,d)$ as a case
where, in their words, ``proving or disproving the conjecture $\dots$ remains a challenge'' ---
every $d\ge3$ falls outside $A_1$-boundedness. We report a machine \emph{verification}, not a
proof, of Conjecture~2 for $249$ positive integral linear forms, including the rusty numbers
$a=(1,d)$ for $d=3,\dots,12$: of the $249$, $61$ --- including every one of these rusty numbers
--- are not $A_1$-bounded and so are reached by no published theorem, while the conjecture in
general is not addressed and remains open. Each
verification guesses a candidate DFAO for the difference sequence and discharges a fixed list
of closed first-order sentences over $\langle\N,+,V_\tau\rangle$ plus the DFAO through a
decision procedure (Peanut) --- $k+9$ of them, eleven for the rusty numbers, where $k=2$; an
induction on $n$ turns these \texttt{TRUE} sentences into a proof that the encoded sequence
\emph{is} the anti-recurrence sequence. One predicate in the first write-up of this method,
\texttt{LAST}, was inverted, silently making two of the obligations vacuous; this was caught in refereeing, the predicate was corrected, and all $249$
records were independently re-verified with the fix. For $a=(1,4)$ we exhibit an explicit
$9$-state base-$5$ DFAO and show it explains, rather than merely reproduces, the anomaly the
published paper reports by hand ($A_{5n+1}=9+55n$ first fails at $n=1741$): the automaton has
exactly one state at which the recursive step deviates, first reached at $n=348$. Two near-miss
automata are exhibited that reproduce the correct sequence for $13\,535\,104$ and $51\,280\,580$
terms respectively and are still wrong, each refuted by an engine-produced counterexample rather
than by longer prefix search --- offered as evidence that the method's positive verifications are
trustworthy. Across the $249$ forms the minimal state count of $a=(1,d)$ obeys
$|Q|(d+3)=|Q|(d)+10$ for every $d\ge3$, an empirical law used to predict, and then confirm,
$|Q|(11)$ and $|Q|(12)$ before those two forms were run.
\end{abstract}

\maketitle

\section{Background and the conjecture}
\label{sec:bg}

For strictly increasing sequences of positive integers $A_n,B_n$ partitioning $\N_{>0}$, and a
positive integral linear form $a=(a_1,\dots,a_k)$, $k>1$, of \emph{trace} $\tau=\sum_j a_j$, the
\emph{anti-recurrence sequence} generated by $a$ is defined by
\[
  A_n \;=\; \sum_{j=1}^{k} a_j\, B_{(n-1)k+j},
\]
where the $n$-th \emph{$B$-block} $B_{(n-1)k+1},\dots,B_{(n-1)k+k}$ is the $k$ smallest
positive integers not yet placed in $A\cup B$ (Fokkink--Joshi's Lemma~3 shows such sequences
exist and are unique, generated blockwise by this mex rule). The anti-Fibonacci numbers,
$a=(1,1)$, are the case popularised by Hofstadter and Zaslavsky \cite{Zaslavsky}; Bosma et al.\
\cite{Bosma} name the general phenomenon the Clergyman's Conjecture and settle $a=(1,1,1)$ and
$a=(1,1,1,1)$. Kimberling and Moses \cite{KimberlingMoses} formulated several questions about
such sequences that together suggest the general meta-conjecture; Fokkink and Joshi \cite{FJ}
give the precise statement:

\begin{conjecture}[{\cite[Conj.~2]{FJ}}]
\label{conj2}
Let $a=(a_1,\dots,a_k)$ be positive and integral of dimension $k>1$. Let $A_n$ be the
anti-recurrence sequence generated by the linear form $f(x)=a\cdot x$. Then $A_n-\kappa n$ is
$\tau$-automatic for $\kappa=k\tau+1$ and $\tau$ the trace of the linear form.
\end{conjecture}

\begin{definition}[{\cite[Def.~1]{FJ}}]
\label{def:a1bdd}
A positive linear form $a$ of dimension $k$ and trace $\tau$ is \emph{$A_1$-bounded} if
$A_1\le(k-1)\tau+2$, where $A_1$ is the first anti-recurrence number in the sequence generated
by $a$.
\end{definition}

\begin{quote}
\textbf{Theorem 4} (Fokkink--Joshi \cite{FJ}). \emph{If $a$ is $A_1$-bounded, then it generates
an anti-recurrence sequence $A_n$ such that $A_n-\kappa n$ is $\tau$-automatic.}
\end{quote}

\noindent The proof of Theorem~4 is constructive: it produces an explicit DFAO recognising the
difference sequence, with at most $2\tau-1$ states.

Section~5 of \cite{FJ} (arXiv~\S4) introduces $a=(1,d)$, ``the rusty numbers,'' exhibits a
guessed $4$-state DFAO for $d=3$, reports that for $d=4$ the subsequence $A_{5n+1}$ agrees with
the arithmetic progression $9+55n$ only up to $n=348$ (``when $A_{1741}\ne9+55\cdot348$''), and
concludes:

\begin{quote}
``Surprisingly, proving or disproving the conjecture for the rusty numbers remains a
challenge.'' \cite[\S5]{FJ}
\end{quote}

\noindent Its closing section asks, verbatim:

\begin{quote}
``Does the conjecture hold without the restriction of $A_1$-boundedness? Are the rusty numbers
sums of linear sequences and automatic sequences?'' \cite[\S6]{FJ}
\end{quote}

For $a=(1,d)$, $k=2$, so $A_1=\sum_j ja_j=1+2d$ and $\tau=1+d$, giving $\kappa=2\tau+1=2d+3$;
$A_1$-boundedness reads $1+2d\le(k-1)\tau+2=(1+d)+2=d+3$, i.e.\ $d\le2$. So $(1,1)$ and $(1,2)$
are covered by Theorem~4 (also, respectively, by Zaslavsky \cite{Zaslavsky} directly and by
Fokkink--Joshi's Theorem~1 for the anti-Pell numbers), and \textbf{every rusty number $d\ge3$
falls outside every published proof of Conjecture~\ref{conj2}} --- exactly the case the
authors call a challenge.

\section{What is established here}
\label{sec:established}

We report machine \emph{verifications} of Conjecture~\ref{conj2} for $249$ positive integral
linear forms $a$, chosen by a sweep over dimension $k=2,\dots,4$ and bounded component size,
including every rusty number $a=(1,d)$, $d=1,\dots,12$. Of these, $61$ forms --- including
$a=(1,d)$ for $d=3,\dots,12$ --- are \emph{not} $A_1$-bounded and so are not covered by any
theorem in \cite{FJ}. ($a=(1,3)$ is a partial exception: \cite{FJ} exhibits a $4$-state DFAO for
it in Figure~7 without a transcript or state count; we independently verify a minimal $6$-state
automaton.)

Each of the $249$ verifications consists of: (i) a candidate DFAO for the difference sequence
$A_n-\kappa n$, produced by guessing the Myhill--Nerode quotient of a computed prefix of the
sequence; (ii) a fixed list of $k+9$ closed first-order sentences over Presburger arithmetic
with a base-$\tau$ numeration predicate ($11$ for the dimension-two forms, including every
rusty number), built from the candidate DFAO, each checked \texttt{TRUE} by the decision
procedure Peanut; and (iii) a hand induction, \S\ref{sec:method} below, showing that
\texttt{TRUE} on those sentences implies the candidate DFAO encodes exactly the mex-generated
pair $(A_n,B_n)$. Steps (ii)--(iii) are what make this a proof for the specific
$a$ verified, rather than a numerical coincidence: unlike prefix matching, a \texttt{TRUE}
first-order sentence over $\langle\N,+,V_\tau\rangle$ decided by the
Bruy\`ere--Hansel--Michaux--Villemaire theorem \cite{BHMV} is a statement about \emph{all}
$n\in\N$, not a finite check.

We want to be exact about what this is and is not. \textbf{It is not a proof of
Conjecture~\ref{conj2}.} What is proved is a large finite list of instances of the conjecture,
one $a$ at a time; nothing here is uniform in $a$, no infinite subfamily is closed, and the
general case --- including whether \emph{every} $a=(1,d)$ works, not just $d\le12$ --- is
exactly as open after this note as before it. What \emph{is} new is that the specific challenge
\cite{FJ} names, the rusty numbers, is machine-verified for the first ten values past the
$A_1$-bounded range, each with an explicit minimal DFAO and a complete, checked induction.

\section{Method: encoding, the corrected predicate, and the induction}
\label{sec:method}

Automatic sequences are $0$-indexed; write $A(m):=A_{m+1}$, $m\ge0$. Fix a base-$\tau$ DFAO $T$
(constants $\mathrm{lo},\mathrm{hi}$; output alphabet $\{0,\dots,\mathrm{hi}-\mathrm{lo}\}$) and
put
\[
  AA(m,s) \;:\equiv\; s=\kappa m+\mathrm{lo}+T[m] \qquad \text{(``}A_{m+1}=s\text{''),}
\]
a finite disjunction over output letters. Write the $n$-th $B$-block as
$x_j=x+(j-1)+[j>g]$, $j=1,\dots,k$, for some $g\in\{1,\dots,k\}$ ($g=k$ meaning no skip,
otherwise the value $x+g$ is skipped); then $A(m)=\tau x+c_g$ with
$c_g=(A_1-\tau)+\sum_{j>g}a_j$. Since $0\le\sum_{j>g}a_j\le\tau-a_1<\tau$ and this sum is
strictly decreasing in $g$, the $c_g$ are pairwise distinct mod~$\tau$, so $(x,g)$ is recovered
from $A(m)$ with no existential quantifier:
\begin{align*}
  G_g(m,x) &:\equiv AA(m,\tau x+c_g) && \text{(one \texttt{let} per } g\text{)}\\
  B1(m,x)  &:\equiv \bigvee_g G_g(m,x) && \text{block start}\\
  LAST(m,y) &:\equiv \bigvee_{g<k}\bigl(\exists x.\,G_g(m,x)\wedge y=x+k\bigr)
              \;\vee\;\bigl(\exists x.\,G_k(m,x)\wedge y=x+k-1\bigr) && \text{block end}\\
  INBLK(m,t) &:\equiv \bigvee_{g,j}\exists u.\,G_g(m,u)\wedge t=u+(j-1)+[j>g]\\
  USED(j,t) &:\equiv AA(j,t)\vee INBLK(j,t)\\
  ISA(t) &:\equiv\exists n.\,AA(n,t) \qquad ISB(t):\equiv\exists m.\,INBLK(m,t)
\end{align*}
Ten fixed closed sentences (functionality/totality of $AA$/$B1$; the base case $G_k(0,1)\wedge
AA(0,A_1)$; monotonicity; $LAST$ strictly below the next block; complementarity of $A$ and $B$;
every value below a block already used) plus one instance of ``the skipped value already used''
per skip-position $g=1,\dots,k-1$ --- $k+9$ sentences in all, eleven when $k=2$, as for every
rusty number --- are then checked \texttt{TRUE}. Standard induction on $m$ --- using $LAST$ to place the new block strictly above
every earlier $B$-value, complementarity to place it strictly above every earlier $A$-value, and
the ``used below $x$'' obligation to force $x=\mex$ --- shows the encoded pair equals the
mex-generated $(A_n,B_n)$ term by term, hence $A_n-\kappa n=(\mathrm{lo}-\kappa)+T[n-1]$ is
$\tau$-automatic: Conjecture~\ref{conj2} for that $a$.

\textbf{The bug, and the fix.} The first write-up of this method, and the script implementing
it, defined $LAST$ by substituting $x:=y+k$ into $G_g$, i.e.\ $LAST(m,y):\equiv\bigvee_{g<k}
G_g(m,y+k)\vee G_k(m,y+k-1)$. That asserts $y=x-k$, the block start \emph{minus} $k$, not the
block's last element $x+k-1$ or $x+k$ --- the wrong direction. The bug made two of the (fixed,
$k$-independent) obligations built from $LAST$ vacuously true for every candidate DFAO, silently
dropping
the step of the induction that uses $LAST$ to place the new block above the previous one; every
other obligation, and the overall \texttt{TRUE}/\texttt{FALSE} verdict, was unaffected in the
sense that no false candidate happened to pass, but the argument as written did not establish
what it claimed to. This was caught in refereeing (\texttt{research/referee-verdicts/attack3-verdict.md}, an
independent, adversarial re-derivation and re-implementation of every step from scratch), which
also confirmed by direct construction that with the inverted predicate the obligation fires on
$0$ of $560$ hand-perturbed near-miss automata, i.e.\ it was carrying no weight. With the
one-line fix above, all $249$ records were rebuilt with an independently written script and
re-run: $249/249$ all obligations \texttt{TRUE}, $0$ failures, total engine time
$541.2$\,s, worst case $14.9$\,s ($a=(1,12)$, $36$ states). The conclusion --- every one of the
$249$ forms verified, including the rusty numbers $d=3,\dots,12$ --- is unchanged by the fix;
what changes is that the argument now actually supports it.

\section{The named case: $a=(1,4)$}
\label{sec:rusty4}

Here $\tau=5$, $\kappa=11$, $A_1=9$; $A_1$-boundedness requires $A_1\le7$, so Theorem~4 does not
apply. A $9$-state base-$5$ DFAO for the difference sequence was guessed and verified (all
eleven obligations \texttt{TRUE} in $0.1$\,s), giving
\[
  A_{m+1} \;=\; 11m+5+T[m].
\]

\textbf{This explains the anomaly \cite{FJ} reports by hand.} Exactly one state $q_8$ of the
nine has $\delta(q_8,0)\ne q_0$ (instead $\delta(q_8,0)=q_1$); every other state maps $0$ back to
$q_0$. Since $\mathrm{state}(5n)=\delta(\mathrm{state}(n),0)$, the subsequence $A_{5n+1}$ is an
arithmetic progression exactly while $\mathrm{state}(n)\ne q_8$, and the first $n$ whose
base-$5$ expansion reaches $q_8$ is $n=348$ (base-$5$ expansion $2343$). Concretely
$A_{1741}=19148=11\cdot1740+5+T[1740]$, one less than $9+55\cdot348=19149$: the automaton
reproduces \cite{FJ}'s reported anomaly to the exact index, and explains it as a single
transition, not a numerical accident.

\textbf{Negative controls.} Two under-sized guesses illustrate why the eleven obligations, and
not prefix agreement, are what makes a verification trustworthy. A $35$-state candidate for
$a=(1,12)$ reproduces the true mex-generated sequence for $51{,}280{,}580$ consecutive terms and
is still wrong: the complementarity obligation and the ``every value below a block already
used'' obligation both return \texttt{FALSE}. Rather than search further by hand, we ask Peanut for a witness to
the negated obligation directly (\texttt{witness \$B1(m,x) \& t>=1 \& t<x \& \~{}(E j. j<m \&
\$USED(j,t))}); it returns the exact breaking index $m=51{,}280{,}580$ in seconds. Re-guessing
from a prefix past that index gives a genuine $36$-state DFAO, which then verifies. A second
example, $a=(2,9)$: a $26$-state candidate reproduces the mex definition for the
$1.2\times10^7$ terms originally computed, and the engine's witness locates its actual first
failure at $m=13{,}535{,}104$; re-guessing past that index gives the correct $27$-state DFAO.
In both cases the engine is used as an equivalence oracle against the mex generator's membership
oracle, so the length of prefix needed is discovered rather than guessed.

\section{Verification}
\label{sec:verif}

Independently of the guess-and-verify pipeline above:

\begin{itemize}
\item A second, separately written mex generator (a naive \texttt{mex}-over-a-\texttt{set},
no pointer arithmetic) agrees with the generator used to produce the $249$ candidates on $A$
and $B$ for the first $3000$ terms of $13$ forms spanning $k=2,3,4$.
\item All $249$ stored DFAOs were rebuilt from their committed \texttt{def} (morphism/coding)
strings and replayed against a fresh run of the naive generator for $n\le20{,}000$: $0$
mismatches, on a code path sharing nothing with the guesser.
\item The generator reproduces published values exactly: OEIS A075326 ($a=(1,1)$), A304502
($a=(1,2)$), A304499 ($a=(2,1)$), A265389 ($a=(1,1,1)$), A299409 ($a=(1,1,1,1)$), and both
published rusty lists, $a=(1,3)$ and $a=(1,4)$, from \cite[\S5]{FJ}.
\item An independent referee re-implementation (\texttt{paper/verdict-attack3/}) built its own
obligation-generating script from scratch, never importing the code above, and its own
generators (a pure-Python \texttt{mex} over a \texttt{set}, and a bitset generator in C); it
re-ran all $249$ records with the corrected \texttt{LAST} and obtained $249/249$
\texttt{TRUE}, and regenerated the $(1,12)$ sequence independently to $7\times10^7$ terms with
$0$ mismatches against the stored $36$-state DFAO.
\end{itemize}

\section{Observations}
\label{sec:obs}

Two structural patterns were noticed across the $249$ verified forms; both are reported as
\textbf{observations on the data}, not as theorems --- neither is proved here, and neither
follows from anything in \S\ref{sec:method}.

\begin{observation}
\label{obs:window}
Across all $249$ forms the difference sequence's output alphabet has width $\mathrm{hi}-
\mathrm{lo}+1\le2\tau-2$, i.e.\ one less than the $2\tau-1$ that Theorem~4 proves under
$A_1$-boundedness; this persists even for the $61$ forms not covered by Theorem~4.
\end{observation}

\begin{observation}[state-count law for the rusty numbers]
\label{obs:law}
The minimal base-$\tau$ DFAO of $a=(1,d)$ has, for $d=1,\dots,12$,
\[
  |Q|(d) \;=\; 2,\,3,\,6,\,9,\,11,\,16,\,19,\,21,\,26,\,29,\,31,\,36 ,
\]
and $|Q|(d+3)=|Q|(d)+10$ for every $d\ge3$ ($7$ instances, $d=3,\dots,9$). The law fails at
$d=1,2$ (the two $A_1$-bounded members: $+7$ and $+8$). \textbf{Two of the seven confirming
instances were predictions made before the corresponding DFAO was computed}: the law, fit to
$d\le9$, predicted $|Q|(11)=31$ and $|Q|(12)=36$; both were then run and confirmed, the second
only after Peanut's counterexample killed the $35$-state near-miss of \S\ref{sec:rusty4}. Every
value is the Myhill--Nerode quotient of a machine-verified automaton, hence an exact minimal
count, not an estimate. The law's own extrapolation gives the untested predictions
$|Q|(13)=39$, $|Q|(14)=41$, $|Q|(15)=46$.
\end{observation}

\section{Limits}
\label{sec:limits}

\textbf{Conjecture~\ref{conj2} is not proved.} What is proved is a finite --- if unusually
large, and individually non-trivial --- list of instances. No argument here is uniform in $a$;
Observation~\ref{obs:law} says a general proof cannot simply relax the constant $2\tau-1$ of
Theorem~4, since the minimal alphabet outgrows it once $A_1$-boundedness fails, so the
substitution alphabet can no longer be read off as ``value of $A_n\bmod\kappa$.'' No infinite
subfamily of the rusty numbers, or of any other family, is closed by this note; whether
\emph{every} rusty number $d$ works remains exactly as open as \cite[\S6]{FJ} leaves it.

Guessing, not verification, is the practical bottleneck as $\tau$ grows: the prefixes needed to
find a correct candidate ranged up to $7\times10^7$ terms ($a=(1,12)$), located by the engine's
counterexample rather than by blind escalation; verifying a found candidate costs seconds to
under a minute. The sweep covers $k=2$ up to $\tau\approx13$, $k=3$ to $\tau=11$, and $k=4$ to
$\tau=9$; larger cases were not attempted, so $249$ is a computational budget boundary, not a
mathematical one. Finally, every claim in this note rests on Peanut being a correct decision
procedure for $\langle\N,+,V_\tau\rangle$ plus a DFAO; that correctness is not formally verified,
though it is supported here by hundreds of adversarial perturbation and random-automaton
controls (all correctly rejected), by exact agreement with an independently written
re-implementation on all $249$ records, and by two independently written sequence generators
agreeing with it everywhere tested.

\subsection*{Acknowledgement and declaration of AI usage}
This note, and the work behind it, were produced with substantial use of a large language
model. This note carries this declaration from its first version. Precisely: the models were Claude (Anthropic): \texttt{claude-fable-5} as the orchestrating model, with \texttt{claude-opus-5} and \texttt{claude-sonnet-5} sub-agents for search, implementation, refereeing and drafting, all used through the Claude Code agent framework; they were used for (i) the literature and prior-art search, including retrieving and checking
the exact published wording of Conjecture~2, Definition~1, and Theorem~4 against the journal
text; (ii) the design and implementation of Peanut, the open-source decision procedure with
which every computation here was run; (iii) the encoding of \S\ref{sec:method}, the
guess-and-verify search for each of the $249$ DFAOs, and the identification of
Observations~\ref{obs:window}--\ref{obs:law}; (iv) an adversarial ``referee'' pass, also carried
out by the model, which re-derived the induction of \S\ref{sec:method} by hand, found the
inverted \texttt{LAST} predicate described there, wrote independent brute-force checks and a
second obligation-generating script sharing no code with the original, and re-ran all $249$
records; and (v) the drafting of this text, most of which is the model's output lightly edited
by the author. The author directed the work, chose the target conjecture, read and takes
responsibility for every statement, and ran the computations. All finite claims were verified by
the independent implementations described in Section~\ref{sec:verif}; the referee transcript and
code are public in the Peanut repository. Any errors are the author's.

\begin{thebibliography}{9}
\setlength{\itemsep}{0pt}\setlength{\parsep}{0pt}\setlength{\parskip}{0pt}

\bibitem{BHMV}
V. Bruy\`ere, G. Hansel, C. Michaux, and R. Villemaire,
\emph{Logic and $p$-recognizable sets of integers},
Bull. Belg. Math. Soc. Simon Stevin \textbf{1} (1994), 191--238.

\bibitem{Bosma}
W. Bosma, R. Bruin, R. Fokkink, J. Grube, A. Reuijl, and T. Tromp,
\emph{Using Walnut to solve problems from the OEIS},
J. Integer Seq. \textbf{28} (2025), Article 25.3.8.

\bibitem{FJ}
R. Fokkink and G. Joshi,
\emph{Anti-recurrence sequences},
Integers \textbf{26} (2026), \#A24;
also arXiv:2506.13337 [math.NT], 16 June 2025.

\bibitem{KimberlingMoses}
C. Kimberling and P. J. C. Moses,
\emph{Linear complementary equations and systems},
Fibonacci Quart. \textbf{57} (2019), no.~5, 97--111.

\bibitem{Zaslavsky}
T. Zaslavsky,
\emph{Anti-Fibonacci numbers, a formula},
unpublished manuscript,
\url{https://oeis.org/A075326/a075326_2.pdf}.

\end{thebibliography}

\end{document}
